\section{Improvement 1: a well-conditioned basis}\label{sec:basis}
\subsection{Basis invariance}
\begin{proposition}[basis invariance]\label{prop:invariance}
If $\bm V$ in~\eqref{eq:vandermonde} is non-singular, the metric functions are the unique
$M_k\in S^e$ with $M_k(\bx_l)=\delta_{kl}$; they, their Lebesgue function $\lambda_S$ and the
element matrix~\eqref{eq:Ke} do not depend on the basis $\{\phi_c\}$.
\end{proposition}
\begin{proof}
For another basis $\bm\psi=\bm C\bm\phi$ with non-singular $\bm C$, $\bm V_\psi=\bm C\bm
V_\phi$ and $\bm\Phi_\psi=\bm C\bm\Phi_\phi$, hence $\bm V_\psi^{-1}\bm\Phi_\psi=\bm
V_\phi^{-1}\bm\Phi_\phi$; the same holds for the gradients.
\end{proof}
Completeness, frame invariance, distortion immunity and unisolvence ($\bm V$ non-singular for
one basis iff for all) are therefore properties of the space $S^e$ and the node set. The basis
affects only the floating-point error, through $\kappa(\bm V)$. Conversely, \emph{no choice of
basis can reduce $\Lambda_S$}: approximation properties and the amplification of nodal
perturbations can only be improved by changing the space (\cref{sec:hybrid}).

\subsection{Candidate bases and scaling}
The original implementation uses monomials $\hat x^a\hat y^b$. For monomials on real nodes
symmetric about the origin, the condition number grows at least exponentially with the degree
\cite{Gautschi1988Lower}; on a parallelogram the 2D matrix is the Kronecker product of two 1D
Vandermonde matrices, so $\kappa_2(\bm V)$ is the square of the 1D value. We consider
tensor-product bases $\phi_a(\tilde s)\phi_b(\tilde t)$ with Legendre polynomials,
Chebyshev polynomials, or the Lagrange polynomials $\ell_a$ through the 1D GLL points. These
polynomials are bounded on $[-1,1]$ but grow rapidly outside, while the skew coordinates of
the nodes of a distorted element leave $[-1,1]^2$ (e.g.\ $|\hat x|\le1.27$ for the trapezoid of
\cref{sec:elements}). The node coordinates are therefore mapped affinely from their bounding
box onto $[-1,1]^2$,
\begin{equation}
  \tilde s=(\hat x-c_s)/r_s ,\qquad \tilde t=(\hat y-c_t)/r_t ,
  \label{eq:bbox}
\end{equation}
with centre $c$ and half-widths $r$ of the box. The map is affine and computed from frame
invariant quantities, so the space $Q_p(\hat x,\hat y)$ and frame invariance are preserved
(\cref{prop:invariance}).

\begin{lemma}[affine elements]\label{lem:affine}
On a parallelogram, the Lagrange basis on the GLL points in the scaled
coordinates~\eqref{eq:bbox} gives $\bm V=\bm\Pi$, a permutation matrix: $\kappa(\bm V)=1$,
$M_k=N_k\circ\bx^{-1}$ exactly, and $\bm K^e$ of~\eqref{eq:Ke} equals the symmetric matrix of the
conventional element.
\end{lemma}
\begin{proof}
For an affine map $\bx=\bx_0+\bm J_0\bxi$, $\hat\bx=\bxi$. The bounding box of the nodes is
$[-1,1]^2$, since the GLL grid contains the end points, so $(\tilde s,\tilde t)=\bxi$ and
$\phi_c(\bx_k)=\ell_a(\xi_k)\ell_b(\eta_k)=\delta_{ck}$ up to the numbering. Then $\bm M=\bm
\Pi^{\T}\bm\Phi$ is evaluated without arithmetic, $M_k(\bx(\bxi))=N_k(\bxi)$, and
$\nabla_{\bx}M_k=\bm J^{-\T}\nabla_{\bxi}N_k$, so $\bm B_L^\ast=\bm B_R\det\bm J$ and
$\bm K^e=\int\bm B_R^{\T}\bm D\bm B_R\det\bm J\dd\xi\dd\eta$.
\end{proof}

\begin{proposition}[distorted elements]\label{prop:kappa}
Let $\hat\bs_d$ be the reference GLL points in the scaled coordinates~\eqref{eq:bbox}, so that
$\phi_c(\hat\bs_d)=\delta_{cd}$, let $\tilde\bs_k$ be the scaled node coordinates, and let
$\lambda_{\mathrm{ref}}(\bs)=\sum_c|\phi_c(\bs)|$. Then
\begin{equation}
  (\bm V^{-1})_{kd}=M_k(\hat\bs_d),\qquad
  \kappa_1(\bm V)=\max_k\lambda_{\mathrm{ref}}(\tilde\bs_k)\;\max_d\lambda_S(\hat\bs_d).
  \label{eq:kappa1}
\end{equation}
\end{proposition}
\begin{proof}
With the query points $\bx_q=\hat\bs_d$, $\bm\Phi=\bm I$ and $\bm M=\bm V^{-1}$. The 1-norm is the
maximum absolute column sum: $\|\bm V\|_1=\max_k\sum_c|\phi_c(\tilde\bs_k)|$ and
$\|\bm V^{-1}\|_1=\max_d\sum_k|M_k(\hat\bs_d)|$.
\end{proof}
Since $\tilde\bs_k\in[-1,1]^2$, the first factor is bounded by the reference value
$\Lambda_{\mathrm{GLL}}$. With this basis, $\kappa(\bm V)$ is therefore the Lebesgue function of
the metric interpolation at the reference points of the bounding box, up to the factor
$\Lambda_{\mathrm{GLL}}$: the basis adds no ill-conditioning beyond that of the interpolation
problem itself. Note that the points $\hat\bs_d$ may lie outside $\Omega^e$, where the metric
polynomials extrapolate, so $\max_d\lambda_S(\hat\bs_d)$ can exceed $\Lambda_S$.

The Legendre and Chebyshev bases have no such identity and are not exact on affine elements
($\kappa_2=55$ and $2.6$ at $p=16$, respectively), but behave similarly on distorted elements
(\cref{sec:elements}). We therefore use the scaled Lagrange basis as the default of the metric
space $Q_p(\hat x,\hat y)$.
